% =====================================================================
%  Thesis-specific packages and macros. Loaded by thesis.tex after the
%  template preamble. Everything here was tested with tagging on.
% =====================================================================
\usepackage{xurl}          % URLs may break at any character
\usepackage{longtable}     % tables that run over a page
\usepackage{array}
% \FloatBarrier keeps a float in its own section. Use the placeins package
% when it is installed; otherwise define the same command here.
\makeatletter
\IfFileExists{placeins.sty}{\usepackage{placeins}}{%
  \def\FloatBarrier{\par\begingroup \let\@elt\relax
    \edef\@tempa{\@botlist\@deferlist\@dbldeferlist}%
    \ifx\@tempa\@empty
    \else
      \ifx\@fltovf\relax
        \if@firstcolumn \clearpage \else \null\newpage\FloatBarrier \fi
      \else
        \newpage \let\@fltovf\relax \FloatBarrier
      \fi
    \fi \endgroup
    \suppressfloats[t]}}
\makeatother

% Set \showlogofalse to drop the university logo from the title page.
% The Library's sample title page has no logo; keep it only if the
% department agrees.
\newif\ifshowlogo
\showlogotrue

% Ragged-right paragraph column of a given width: L{4cm}
\newcolumntype{L}[1]{>{\raggedright\arraybackslash}p{#1}}

% Identifiers, commands and short file names. Underscores need no
% escape: \code{sf_bench}. Use \nolinkurl{...} for long paths.
\NewDocumentCommand{\code}{m}{\texttt{\detokenize{#1}}}

% Microseconds.
\newcommand{\us}{\,\textmu s}

\emergencystretch=2em
\hyphenation{lib-fabric lib-ib-verbs Soft-RoCE Infini-Band Flash-Atten-tion}

% Tables: a little air between rows, so that cells of several lines do
% not run into the next row; captions of long tables as wide as the text;
% and space between a caption and the table under it.
\AddToHook{env/tabular/begin}{\renewcommand{\arraystretch}{1.3}}
\AddToHook{env/longtable/begin}{\renewcommand{\arraystretch}{1.3}}
\setlength{\LTcapwidth}{\linewidth}
\AddToHook{env/table/begin}{\setlength{\abovecaptionskip}{2pt}\setlength{\belowcaptionskip}{8pt}}

% Run-in labels. \paragraph under a \section would be a level-5 heading
% directly below a level-2 one, which skips heading levels in the tagged
% PDF, so a label is set as italic text at the start of the paragraph
% instead of as a heading. The text of the thesis does not use run-in
% labels or bold lead-ins: a paragraph opens with a topic sentence. Inside
% a proof the steps are named in italics with \emph{Step name.}
\RenewDocumentCommand{\paragraph}{s o m}{%
  \par\noindent\textit{#3}\hspace{0.6em}\ignorespaces}

% The lemma and definition environments are declared in thesis.tex, next
% to the proposition environment.

% --- Headings in the form of the department's theses ----------------------
% A chapter opens with CHAPTER and its number on one centred line and the
% title centred under it, in bold at a size close to that of the text.
% Section headings are bold and flush left. The names of the unnumbered
% parts (contents, lists, bibliography) are set in capitals.
% TeX Live 2026 builds headings from templates and TeX Live 2025 from the
% macros of the report class, so both are set. Neither touches the tags.
\renewcommand{\contentsname}{CONTENTS}
\renewcommand{\listfigurename}{LIST OF FIGURES}
\renewcommand{\listtablename}{LIST OF TABLES}
\renewcommand{\bibname}{BIBLIOGRAPHY}
\makeatletter
\IfInstanceExistsTF{heading}{chapter}
  {%
    \EditInstance{heading}{chapter}{
      after-penalty-vspace = 36pt,
      after-vspace         = 30pt,
      text                 = \MakeUppercase{\@chapapp},
      heading-decls        = \centering \parindent0pt \normalfont \bfseries \large,
      title-decls          = \large,
    }
    \EditInstance{headformat}{chapter}{
      separator-b = \par \nobreak \vspace{2pt},
    }
    \EditInstance{heading}{part}{
      text          = \MakeUppercase{\partname},
      heading-decls = \centering \normalfont \bfseries \Large,
      title-decls   = \Large,
    }
    \EditInstance{heading}{section}{heading-decls = \normalfont \large \bfseries}
    \EditInstance{heading}{subsection}{heading-decls = \normalfont \normalsize \bfseries}
    \EditInstance{heading}{subsubsection}{heading-decls = \normalfont \normalsize \itshape}
  }
  {%
    \def\@makechapterhead#1{%
      \vspace*{24\p@}%
      {\parindent \z@ \centering \normalfont
        \ifnum \c@secnumdepth >\m@ne
          \large\bfseries \MakeUppercase{\@chapapp}\space \thechapter
          \par\nobreak
          \vskip 2\p@
        \fi
        \interlinepenalty\@M
        \large \bfseries #1\par\nobreak
        \vskip 30\p@
      }}
    \def\@makeschapterhead#1{%
      \vspace*{24\p@}%
      {\parindent \z@ \centering \normalfont
        \interlinepenalty\@M
        \large \bfseries #1\par\nobreak
        \vskip 30\p@
      }}
    \def\@part[#1]#2{%
      \ifnum \c@secnumdepth >-2\relax
        \refstepcounter{part}%
        \addcontentsline{toc}{part}{\thepart\hspace{1em}#1}%
      \else
        \addcontentsline{toc}{part}{#1}%
      \fi
      \markboth{}{}%
      {\centering
       \interlinepenalty \@M
       \normalfont
       \ifnum \c@secnumdepth >-2\relax
         \Large\bfseries \MakeUppercase{\partname}\nobreakspace\thepart
         \par
         \vskip 30\p@
       \fi
       \Large \bfseries #2\par}%
      \@endpart}
    \renewcommand\section{\@startsection{section}{1}{\z@}%
      {-3.5ex \@plus -1ex \@minus -.2ex}{2.3ex \@plus .2ex}%
      {\normalfont\large\bfseries}}
    \renewcommand\subsection{\@startsection{subsection}{2}{\z@}%
      {-3.25ex \@plus -1ex \@minus -.2ex}{1.5ex \@plus .2ex}%
      {\normalfont\normalsize\bfseries}}
    \renewcommand\subsubsection{\@startsection{subsubsection}{3}{\z@}%
      {-3.25ex \@plus -1ex \@minus -.2ex}{1.5ex \@plus .2ex}%
      {\normalfont\normalsize\itshape}}
  }
\makeatother

% --- The research questions -------------------------------------------------
% Stated once, so that the problem statement and the conclusion repeat them
% in the same words.
\newcommand{\rqone}{Does giving each traffic class its own endpoint keep the
  tail latency of small messages flat while bulk transfers run, on software
  RoCE and on hardware RoCE?}
\newcommand{\rqtwo}{What does a portable transport interface cost against the
  raw RDMA interface on the path small messages take?}
\newcommand{\rqthree}{How much of the time to first token is cache transfer,
  and do per-transfer routing, compression and spilling reduce it against a
  transport that treats every transfer alike?}

% --- Star notes: things to fix, fill in or remove ------------------------
% Every note printed in the PDF starts with a star (*), and nothing else in
% the document uses that character, so a search for * finds everything that
% still has to be fixed, filled in or removed. Three commands print notes:
%   \fixnote{text}             * text
%   \waiting{results from E3}  * Waiting for results from E3.
%   \pendingfigure{...}        a placeholder figure whose box starts with
%                              "* Waiting for ..." and lists the steps
% \blank prints a blank for a value that a run will supply; the sentence is
% followed by \waiting{...}. Every note is also written to thesis.aux, and
%   python3 tools/notes.py thesis.aux
% lists them with their page numbers. To remove a note, delete its command.
% \shownotesfalse hides the notes in the text and under the tables, which is
% only for checking the page count: the blanks, the empty cells and the
% placeholder figures stay.
\newif\ifshownotes
\shownotestrue
\newcommand{\blank}{\mbox{\_\_\_\_\_}}
\makeatletter
\newcommand{\thesisnote}[2]{}   % read back from the .aux file; prints nothing
\NewDocumentCommand{\thesislognote}{+m}{%
  \protected@write\@auxout{}{\string\thesisnote{\thepage}{\detokenize{#1}}}}
\makeatother
\NewDocumentCommand{\fixnote}{+m}{%
  \ifshownotes
    \thesislognote{#1}%
    \textbf{*}\,\textit{#1}%
  \fi}
\NewDocumentCommand{\waiting}{m}{\fixnote{Waiting for #1.}}

% --- Placeholder for a figure that waits for a result --------------------
% \pendingfigure{label}{caption}{what it waits for}{step one ;; step two ;; ...}
% Draws a dashed box that starts with "* Waiting for <what it waits for>."
% and lists the steps that produce the real figure. The same words become
% the alt text. The third argument completes the sentence "Waiting for ...",
% for example: results from E3 (run R2). The last two arguments are read
% verbatim: plain text only, underscores allowed, steps separated by two
% semicolons. Use it at the top level of a chapter, never inside another
% command.
\newcommand{\thesispendingbox}[3]{%
  \begin{tikzpicture}[diagram, alt={#1}]
    \node[draw=dInk, dashed, line width=0.8pt, fill=white,
          text width=0.88\linewidth, align=left, inner sep=9pt, font=\small]
      {\textbf{*\,Waiting for #2.}\par\smallskip
       To produce this figure:\par #3};
  \end{tikzpicture}}
\ExplSyntaxOn
\seq_new:N \l__thesis_steps_seq
\tl_new:N \l__thesis_alt_tl
\tl_new:N \l__thesis_disp_tl
\int_new:N \l__thesis_step_int
\cs_new_protected:Npn \__thesis_box:nnn #1#2#3 { \thesispendingbox {#1} {#2} {#3} }
\cs_generate_variant:Nn \__thesis_box:nnn { VnV }
\NewDocumentCommand{\pendingfigure}{m m +v +v}
  {
    \seq_set_split:Nnn \l__thesis_steps_seq { ;; } {#4}
    \tl_set:Nn \l__thesis_alt_tl { Placeholder~for~a~figure.~Waiting~for~#3.~To~produce~it:~ }
    \tl_clear:N \l__thesis_disp_tl
    \int_zero:N \l__thesis_step_int
    \seq_map_inline:Nn \l__thesis_steps_seq
      {
        \int_incr:N \l__thesis_step_int
        \tl_put_right:Ne \l__thesis_alt_tl { Step~\int_use:N \l__thesis_step_int :~ \exp_not:n {##1} .~ }
        \tl_put_right:Ne \l__thesis_disp_tl { \int_use:N \l__thesis_step_int .~ \exp_not:n {##1} \exp_not:N \par }
      }
    % two hyphens in a command-line option must not join into a dash
    \tl_replace_all:Nnn \l__thesis_disp_tl { -- } { - \kern 0pt - }
    \begin{figure}[htbp]
      \centering
      \thesislognote{Waiting~for~#3.~(placeholder~figure:~#2)}
      \__thesis_box:VnV \l__thesis_alt_tl {#3} \l__thesis_disp_tl
      \caption[#2]{#2.}
      \label{#1}
    \end{figure}
  }
\ExplSyntaxOff

% Single-spaced, small block for long tables outside floats.
\newenvironment{tightblock}{\par\begingroup\setstretch{1}\small}{\par\endgroup}

% Reference names, written out and capitalised.
\crefname{figure}{Figure}{Figures}
\crefname{table}{Table}{Tables}
\crefname{chapter}{Chapter}{Chapters}
\crefname{section}{Section}{Sections}
\crefname{appendix}{Appendix}{Appendices}
\crefname{equation}{Equation}{Equations}
\crefname{proposition}{Proposition}{Propositions}
\crefname{definition}{Definition}{Definitions}
\crefname{lemma}{Lemma}{Lemmas}
\crefname{algorithm}{Algorithm}{Algorithms}
\crefname{part}{Part}{Parts}

% Plot data generated from the results files in the repository.
% Generated from the results files in the repository. Do not edit by hand.
\newcommand{\quantOverlapCoords}{(0,90.62) (1,82.81) (2,90.62) (3,85.94) (4,96.88) (5,92.19) (6,90.62) (7,90.62) (8,89.06) (9,93.75) (10,89.06) (11,87.50) (12,89.06) (13,89.06) (14,93.75) (15,89.06) (16,92.19) (17,89.06) (18,92.19) (19,90.62) (20,96.88) (21,92.19) (22,87.50) (23,87.50) (24,93.75) (25,87.50) (26,87.50) (27,92.19) (28,90.62) (29,89.06) (30,87.50) (31,92.19) (32,89.06) (33,96.88) (34,93.75) (35,93.75) (36,95.31) (37,93.75) (38,96.88) (39,95.31)}
\newcommand{\quantErrSketchCoords}{(0,25.29) (1,80.80) (2,31.94) (3,28.18) (4,21.61) (5,24.54) (6,17.55) (7,17.07) (8,23.40) (9,20.31) (10,18.78) (11,23.53) (12,24.58) (13,26.76) (14,25.55) (15,28.12) (16,24.23) (17,27.97) (18,27.49) (19,24.19) (20,27.07) (21,27.25) (22,20.98) (23,23.74) (24,24.07) (25,24.31) (26,20.97) (27,21.14) (28,21.06) (29,26.20) (30,24.65) (31,21.28) (32,19.77) (33,23.28) (34,19.23) (35,17.08) (36,23.07) (37,24.84) (38,23.07) (39,29.23)}
\newcommand{\quantErrPlainCoords}{(0,65.39) (1,205.51) (2,71.05) (3,56.68) (4,47.07) (5,52.73) (6,38.24) (7,41.20) (8,56.98) (9,48.39) (10,41.93) (11,61.35) (12,47.31) (13,59.01) (14,62.77) (15,66.43) (16,69.66) (17,67.76) (18,72.02) (19,63.13) (20,59.76) (21,53.74) (22,41.50) (23,53.97) (24,52.23) (25,51.42) (26,45.08) (27,46.67) (28,47.55) (29,52.22) (30,47.80) (31,53.25) (32,42.93) (33,63.07) (34,44.50) (35,41.40) (36,52.02) (37,50.77) (38,42.11) (39,91.43)}
\newcommand{\nvmeWriteCoords}{(1,662.1) (10,730.7) (20,802.7) (50,827.4)}
\newcommand{\nvmeReadCoords}{(1,1731.3) (10,2149.4) (20,2187.6) (50,1103.9)}
\newcommand{\decompGpuSone}{(4096,0.459) (32768,2.407) (131072,20.763)}
\newcommand{\decompGpuStwo}{(4096,1.468) (32768,3.977) (131072,33.583)}
\newcommand{\decompCpuSone}{(4096,6.641) (32768,56.478) (131072,572.297)}
\newcommand{\decompCpuStwo}{(4096,8.928) (32768,135.442) (131072,1026.065)}


% \plotcoords[plot options]{\macroHoldingCoordinates}: an \addplot whose
% coordinate list is stored in a macro (see figures/data.tex).
\makeatletter
\newcommand{\plotcoords}[2][]{\expandafter\thesis@plot\expandafter{#2}{#1}}
\newcommand{\thesis@plot}[2]{\addplot[#2] coordinates {#1};}
\makeatother

% Shared look of the charts.
\pgfplotsset{thesis chart/.style={
    width=0.84\linewidth, height=6.0cm,
    axis lines*=left, axis line style={draw=dInk, line width=0.7pt},
    ymajorgrids, grid style={draw=dGray!35, line width=0.4pt},
    tick align=outside, tick style={draw=dInk},
    tick label style={font=\footnotesize}, label style={font=\small},
    legend style={draw=none, fill=white, font=\footnotesize, legend cell align=left, row sep=1pt, inner sep=3pt},
  }}
